% Lab 4 — Constraining Multiple Clocks（原书 PDF p.51–70）
\bichapter{Lab 4：约束多时钟}{Lab 4: Constraining Multiple Clocks}

\begin{objectiveBox}
完成本实验后，你应能够：
\begin{itemize}
  \item 运用课堂讲授的命令收集设计中时钟的相关信息；
  \item 使用 GUI 从另一个视角查看设计时钟及其关系；
  \item 识别并报告虚假违例（false violation）；
  \item 确认分析结果。
\end{itemize}
\end{objectiveBox}

\begin{durationBox}
约 \textbf{45} 分钟。
\end{durationBox}

\section{概述}
\label{lab4:overview}

\noindent\textit{Overview}

本实验大致步骤如下：
\begin{enumerate}
  \item 在 shell 中熟悉设计时钟；
  \item 在 GUI 中查看时钟关系；
  \item 使用 GUI 深入探索时序路径细节；
  \item 报告虚假违例；
  \item 重新执行运行脚本以减少违例。
\end{enumerate}

\subsection{相关文件与目录}
\label{lab4:files}

\noindent\textit{Relevant Files and Directories}

本实验全部文件位于主目录下的 \cmd{lab4_clocks} 目录中。

\begin{tabularx}{\textwidth}{@{}l X@{}}
\toprule
\textbf{路径 / 文件} & \textbf{说明} \\
\midrule
\cmd{lab4_clocks/} & 当前工作目录 \\
\cmd{orca_savesession/} & ORCA 初始保存会话 \\
\cmd{orca_savesession_violations/} & 含问题的 ORCA 保存会话 \\
\cmd{RUN.tcl} & ORCA 运行脚本 \\
\cmd{.synopsys\_pt.setup} & PT setup 文件 \\
\cmd{scripts/} & \\
\quad\cmd{orca\_pt\_variables.tcl} & 变量脚本 \\
\bottomrule
\end{tabularx}


\section{操作说明}
\label{lab4:instructions}

\noindent\textit{Instructions}

\subsection{任务 1：熟悉设计时钟}
\label{lab4:task1}

\noindent\textit{Task 1. Get to Know the Design Clocks}

\begin{enumerate}
  \item 确认当前目录为 \cmd{lab4_clocks}。

  \item 启动 PrimeTime（\cmd{pt\_shell}）。
        恢复保存在 \cmd{./orca\_savesession} 中的会话。
        可键入少量字符后按 \textbf{Tab} 键利用命令与文件名补全。

  \item 运用课堂讲授的命令回答下列问题。
        可借助 Job Aid「Clocks and More」回忆具体命令。

\begin{questionBox}
\textbf{问题 1.}\ 设计中有多少个时钟？其中有多少个是 generated clock？\\[0.35em]
\textbf{答：} 可用 \cmd{report\_clock} 或下列命令获取：
\begin{lstlisting}
pt_shell> sizeof_collection [all_clocks]
6
pt_shell> sizeof_collection [get_generated_clocks *]
3
\end{lstlisting}
设计中共有 \textbf{6} 个时钟，其中 \textbf{3} 个为 generated clock。

\textbf{问题 2.}\ 哪些输入端口定义了 master clock？\\[0.35em]
\textbf{答：}\\
\begin{lstlisting}
pt_shell> rpt_clock_ports
Port Name  Direction  Clock Name    Is Generated
pclk       in         PCI_CLK       false
sys_clk    in         SYS_CLK       false
sdr_clk    in         SDRAM_CLK     false
sd_CK      out        SD_DDR_CLK    true
sd_CKn     out        SD_DDR_CLKn   true
\end{lstlisting}
定义了 master clock 的输入端口为 \textbf{pclk}、\textbf{sys\_clk}、\textbf{sdr\_clk}。

\textbf{问题 3.}\ 哪些输出端口定义了 outgoing clock？\\[0.35em]
\textbf{答：} 由同一报告可知，输出端口为 \textbf{sd\_CK} 与 \textbf{sd\_CKn}。

\textbf{问题 4.}\ 时钟是 propagated 还是 ideal？\\[0.35em]
\textbf{答：} 用 \cmd{report\_clock} 可见，设计中所有时钟均为 \textbf{ideal}。

\textbf{问题 5.}\ 哪 3 对时钟之间存在已约束的时序路径？\\[0.35em]
\textbf{答：}\\
\begin{lstlisting}
pt_shell> check_timing -override clock_crossing -verbose
Information: Checking 'clock_crossing'.
Information: There are 4 clocks having domains interacting.
# all paths are false paths
# part of paths are false paths
From Clock    Crossing Clocks
PCI_CLK       SYS_CLK*
SDRAM_CLK     SD_DDR_CLK#, SYS_CLK*
SYS_2x_CLK    SDRAM_CLK*, SYS_CLK
SYS_CLK       PCI_CLK*, SDRAM_CLK*, SYS_2x_CLK
\end{lstlisting}
存在已约束时序路径的 3 对时钟为：
\textbf{SDRAM\_CLK} 与 \textbf{SYS\_CLK}；
\textbf{SYS\_2x\_CLK} 与 \textbf{SYS\_CLK}；
\textbf{SYS\_CLK} 与 \textbf{PCI\_CLK}（以及 \textbf{SDRAM\_CLK}、\textbf{SYS\_2x\_CLK}）。
\end{questionBox}
\end{enumerate}

\subsection{任务 2：使用 GUI 报告时钟关系}
\label{lab4:task2}

\noindent\textit{Task 2. Use the GUI to Report Clock Relationships}

若设计中时钟较多，GUI 可简化理解时钟之间关系的任务。

\begin{enumerate}
  \item 执行下列命令启动 GUI：
\begin{lstlisting}
pt_shell> start_gui
\end{lstlisting}

\begin{noteBox}
原始 \cmd{pt\_shell} 会话仍在终端窗口中运行。
可保持 GUI 打开，并根据任务需要选择 shell 或 GUI 界面。
\end{noteBox}

  \item 查看时钟域交叉（clock domain crossings）：
        从下拉菜单打开 \textbf{Clock > Clock Analyzer} 中的「clock domain matrix」。
        打开的 Clock Analyzer 窗口（必要时点击时钟左侧的加号展开）应与
        \cmd{check\_timing} 报告设计内时钟交叉时的信息一致。
        将鼠标悬停在矩阵方块上可查看存在何种 false path。
        与 \cmd{check\_timing -override clock\_crossing -verbose} 的文本输出相比，
        图形化矩阵有时更易于理解。

        窗口左侧列出每个 master clock 以及由该 master clock 派生的 generated clock。

\begin{questionBox}
\textbf{问题 6.}\ \cmd{SYS\_2x\_CLK} 的 master clock 是什么？\\[0.35em]
\textbf{答：} \textbf{SYS\_CLK}。

\textbf{问题 7.}\ \cmd{SYS\_2x\_CLK} 定义在哪个 pin/port（其「source」）上？
（提示：可能需要拖开时钟矩阵窗口以露出更多时钟信息列。）\\[0.35em]
\textbf{I\_CLOCK\_GEN/I\_CLKMUL/CLK\_2x}。
（可能需要拖开矩阵窗口才能看到 source pin。）

\textbf{问题 8.}\ \cmd{SYS\_2x\_CLK} 的 master clock 定义在哪个 pin/port 上？\\[0.35em]
\textbf{答：} \textbf{sys\_clk}。
\end{questionBox}

  \item 进一步探索：显示 \cmd{SYS\_2x\_CLK} 的时钟原理图——
        选中该时钟，然后右键 \textbf{Schematic of Selected Clocks}。
        展开名为 \cmd{I\_CLOCK\_GEN/U20} 的 MUX 的 fanin
        （提示：可用 \textbf{Select > By Name} 定位/高亮 U20）：
        双击输入 stub，如下图所示，继续双击直至 fanin 耗尽（例如已到达输入端口）。

\begin{noteBox}
图例说明：\texttt{[ce]} 表示 fully constrained paths；
\texttt{Launch:} 表示启动时钟；矩阵中各时钟块显示周期与波形信息。
在原理图中双击输入 pin stub 可继续向上展开 fanin。
\end{noteBox}

\begin{questionBox}
\textbf{问题 9.}\ 哪个端口连接到 MUX \cmd{I\_CLOCK\_GEN/U20} 的 select pin？\\[0.35em]
\textbf{答：} \textbf{test\_mode}。

\textbf{问题 10.}\ 查看原理图是否有助于理解测试模式下的时钟方案？\\[0.35em]
\textbf{答：} 是——两个时钟进入 MUX：一个来自时钟发生器，一个直接来自端口。
\cmd{test\_mode} 端口控制 select 线，使测试模式下可旁路时钟发生器，
让设计直接由端口驱动。
\end{questionBox}

  \item 用抽象时钟图（abstract clock graph）探索时钟关系：
        关闭原理图窗口，在 TopLevel 窗口中选择 \textbf{Clock > Clock Graph for All Clocks}。
        必要时按 \textbf{F8} 显示工具栏，勾选工具栏中的各项后点击 \textbf{Apply}。

  \item 在 Abstract Clock Graph 工具栏中选择 \textbf{Mux} 并点击 \textbf{Apply}，找出一对 muxed clock。

  \item 在 Abstract Clock Graph 中定位实例 \cmd{I\_CLOCK\_GEN/U10}（\cmd{mx02d1}）。
        （提示：可用 \textbf{Select > By Name} 定位/高亮 U10。）

\begin{questionBox}
\textbf{问题 11.}\ 哪些时钟驱动 \cmd{I\_CLOCK\_GEN/U10}？\\[0.35em]
\textbf{答：} \textbf{SYS\_CLK} 与 \textbf{SYS\_2x\_CLK}。
\end{questionBox}

  \item 从时钟图窗口「放大」到感兴趣的对象：
        选中 \cmd{I\_CLOCK\_GEN/U10}，然后 \textbf{Schematic > Schematic View}。

\begin{questionBox}
\textbf{问题 12.}\ 哪个端口驱动 \cmd{I\_CLOCK\_GEN/U10} 的 select 线？\\[0.35em]
\textbf{答：} \textbf{power\_save}。
\end{questionBox}

  \item 返回 Abstract Clock Graph。

\begin{questionBox}
\textbf{问题 13.}\ 从 abstract clock graph 窗口，能否打开并显示与 Clock Analyzer 中相同的
\cmd{SYS\_2x\_CLK} 时钟原理图？（在 \cmd{SYS\_2x\_CLK} 上右键查找相应选项。）\\[0.35em]
\textbf{答：} 可以。选中 \cmd{SYS\_2x\_CLK}（可能需要放大以仅选中该时钟），
右键选择 \textbf{Schematic for Selected Clocks}。
可能需要双击展开输入或输出 stub 才能得到与 Clock Analyzer 中完全相同的原理图。
\end{questionBox}

  \item 点击 Clock Analyzer 窗口右上角的 ``X'' 关闭该窗口。

  \item 点击 Clock Schematic 与 Clock Analyzer 窗口右上角的 ``X'' 关闭这些窗口。
\end{enumerate}

\subsection{任务 3：使用 GUI 探索时序路径细节}
\label{lab4:task3}

\noindent\textit{Task 3. Use the GUI to Explore Detail of Timing Paths}

研究启动时钟与捕获时钟之间的路径——本任务将查看由 \cmd{SYS\_CLK} 时钟的启动与捕获路径上的网络延迟（network latency）。

\begin{enumerate}
  \item 在 shell 中（GUI 背后的终端窗口仍保持打开）执行下列命令，
        在检查路径之前将所有时钟设为 propagated，以便 PrimeTime 计算时钟网络延迟
        （约需一分钟）。
        令 PrimeTime 保存所有 pin 的 arrival time（这正是你将要检查的内容），
        然后定义待检查的时序路径集合：
\begin{lstlisting}
set_propagated_clock [all_clocks]
set timing_save_pin_arrival_and_slack true
update_timing
set my_paths [get_timing_paths -max 10 -group SYS_CLK \
    -path full_clock_expanded]
\end{lstlisting}

  \item 从下拉菜单 \textbf{Timing > Path Analyzer} 输入你的违例路径集合。
        输入时序路径集合并点击 \textbf{Apply}。

  \item 为这 10 条时序路径生成直方图：
        在 Path Analyzer 中右键点击 \textbf{ALL}，选择 \textbf{Create Histogram}，
        然后 \textbf{Apply}；在弹出的 Categories Table/Histogram 对话框中点击 \textbf{OK}。

  \item 从直方图中对选定的路径启动 Path Inspector：
        左键点击最左侧（最差 slack）的柱条以选中该组路径；
        再选中最差路径，点击 \textbf{Inspector}。

  \item 在 Path Inspector 中检查时钟重汇聚悲观（clock reconvergent pessimism, CRP）：
        在 data required 与 data arrival 段中向下滚动找到 CRP，
        再横向滚动查看 CRP 占延迟的百分比。

\begin{questionBox}
\textbf{问题 14.}\ capture delay 中有百分之多少来自 CRP？\\[0.35em]
\textbf{答：} \textbf{4.18\%}。

\textbf{问题 15.}\ 该百分比是否对所有设计都具有代表性？\\[0.35em]
\textbf{答：} 否——该数值取决于具体设计与具体路径。
\end{questionBox}

  \item 点击 Path Inspector 窗口底部的 \textbf{Schematic} 标签查看路径原理图。

  \item 在原理图窗口中找到 CRP 点——即启动路径与捕获路径分叉前的最后一个 pin。

\begin{noteBox}
鼠标手势（gestures/strokes）可简化缩放：
按住中键并垂直拖动光标可「全屏缩放」；
沿对象对角向上拖动可放大，向下拖动可缩小。
\end{noteBox}

  \item 要查看该 pin 上的 arrival time，必要时须先「展开」该 pin 的 buffer
        （默认 PrimeTime 会将 buffer 树折叠为单个符号）。

  \item 选中寄存器前一级 buffer 的输出 pin，然后选择 \textbf{View > Property}，
        查看 arrival time 及其他感兴趣的属性。

\begin{questionBox}
\textbf{问题 16.}\ 该 buffer 输出 pin 的 arrival window 有多宽？\\[0.35em]
\textbf{答：} 为 $2.86939 - 2.40107 \approx \mathbf{0.468}$\,ns。

\textbf{问题 17.}\ 这是否与先前在 Path Inspector 的 data arrival / data required 段中看到的一致？\\[0.35em]
\textbf{答：} 是。
\end{questionBox}

  \item 检查路径波形：点击 Path Inspector 窗口底部的 \textbf{Waveform} 标签。

\begin{questionBox}
\textbf{问题 18.}\ 在波形窗口中右键点击可添加什么内容？\\[0.35em]
\textbf{答：} 可添加输入 pin 与输出 pin，从而查看整条路径的细节。
\end{questionBox}

  \item 关闭 GUI 但保持原始 \cmd{pt\_shell} 会话在终端中继续运行：
\begin{lstlisting}
pt_shell> stop_gui
\end{lstlisting}
        或在主 GUI 窗口选择 \textbf{File > Close GUI}。

  \item 退出 PrimeTime。
\end{enumerate}

\subsection{任务 4：报告虚假违例}
\label{lab4:task4}

\noindent\textit{Task 4. Report a False Violation}

\begin{enumerate}
  \item 启动 PrimeTime 并恢复保存在 \cmd{orca\_savesession\_violations} 中的会话。

  \item 确定 ORCA 中时序违例的数量与类型：
\begin{lstlisting}
pt_shell> report_analysis_coverage
\end{lstlisting}

\begin{questionBox}
\textbf{问题 19.}\ ORCA 有多少个违例？分别是什么类型？\\[0.35em]
\textbf{答：} 共有 \textbf{23} 个 setup 违例、\textbf{53} 个 hold 违例、\textbf{32} 个 out\_setup 违例。
\end{questionBox}

  \item 为 \cmd{out\_setup} 时序检查中最差 slack 生成一份「简短（short）」时序报告。

\begin{questionBox}
\textbf{问题 20.}\ 如何找出 \cmd{out\_setup} 中端点 port 的最差 slack？
（可借助 Job Aid「Timing Reports」回忆两个合适的开关。）\\[0.35em]
\textbf{答：}\\
\begin{lstlisting}
pt_shell> report_timing -path short -to [all_outputs]
\end{lstlisting}
或：
\begin{lstlisting}
pt_shell> page_on
pt_shell> report_analysis_coverage -status violated -check out_setup
Type of Check  Total  Met        Violated
out_setup      75     43 (57%)   32 (43%)
\end{lstlisting}

\textbf{问题 21.}\ 该违例涉及哪些时钟（launch 与 capture）？\\[0.35em]
\textbf{答：}\\
\begin{lstlisting}
pt_shell> report_timing -to sd_DQ[6] -path short
Startpoint: sdr_clk (clock source 'SDRAM_CLK')
Endpoint:   sd_DQ[6] (output port clocked by SD_DDR_CLK)
Path Group: SD_DDR_CLK
Path Type:  max
\end{lstlisting}
Launch 时钟为 \textbf{SDRAM\_CLK}，capture 时钟为 \textbf{SD\_DDR\_CLK}。
\end{questionBox}

\begin{noteBox}
由任务 1 可知，\cmd{SD\_DDR\_CLK} 是定义在输出端口上的 generated clock。
定义 outgoing clock 的目的是让 PrimeTime 计算该时钟的 source latency，
并将其纳入 data required time。
\end{noteBox}

  \item 查看上一步时序报告的 data required time 段，注意其中未报告 clock latency。
        用下列命令确认：
\begin{lstlisting}
# 该报告将无输出，因为 PrimeTime 尚未为 SD_DDR_CLK 计算 source latency
pt_shell> report_clock -skew SD_DDR_CLK
\end{lstlisting}

\begin{questionBox}
\textbf{问题 22.}\ 为何 PrimeTime 未为 outgoing clock \cmd{SD\_DDR\_CLK} 计算 source latency？\\[0.35em]
\textbf{答：} 必须为 propagated clock（尤其是 master clock），
PrimeTime 才能为 generated clock 计算 source latency。
本设计中所有时钟均为 ideal。
\end{questionBox}

\begin{noteBox}
与设计师沟通后发现存在误解：设计师期望你打开一个变量，使所有时钟均为 propagated！
\end{noteBox}

  \item 有一个变量可使所有时钟变为 propagated。用 Tcl 过程 \cmd{aa} 帮助定位该变量：
\begin{lstlisting}
pt_shell> aa propagate
\end{lstlisting}

\begin{questionBox}
\textbf{问题 23.}\ 该变量的名称是什么？\\[0.35em]
\textbf{答：} \textbf{timing\_all\_clocks\_propagated}。
\begin{lstlisting}
pt_shell> aa propagate
...
timing_all_clocks_propagated = "false"
\end{lstlisting}

\textbf{问题 24.}\ 根据 man page，该变量将做什么？\\[0.35em]
\textbf{答：} 将该变量设为 \texttt{true} 后，之后创建的所有时钟都将以 propagated clock 方式创建。
（下一任务中将实际设置该变量。）
\end{questionBox}

  \item 查阅 \cmd{check\_timing} 的 man page，找出会标记所有 ideal clock 的附加检查名称。
        下列命令在带滚动条的弹出窗口中打开 man page，便于阅读长文档：
\begin{lstlisting}
pt_shell> vman check_timing
\end{lstlisting}

\begin{noteBox}
上述命令为 \cmd{.synopsys\_pt.setup} 中定义的别名，使用 SolvNet Doc Id 014947 上的 \cmd{view} 命令。
若环境中未安装或未提供 Tk 包的主可执行文件 \cmd{wish}，别名 \cmd{vman} 将不可用。
\end{noteBox}

\begin{questionBox}
\textbf{问题 25.}\ 如何修改 \cmd{check\_timing} 以添加检查，验证所有时钟均为 propagated？\\[0.35em]
\textbf{答：} 附加检查名为 \textbf{ideal\_clocks}。
可用 \cmd{lappend} 将其加入变量 \cmd{timing\_check\_defaults}，
使 \cmd{check\_timing} 自动执行该检查（下一任务中将在
\cmd{./scripts/orca\_pt\_variables.tcl} 中完成）：
\begin{lstlisting}
lappend timing_check_defaults ideal_clocks
set timing_all_clocks_propagated true
\end{lstlisting}

\textbf{任务 5 步骤 4 验证要点：}
\begin{itemize}
  \item \cmd{out\_setup} 违例应已减少；
  \item \cmd{report\_clock} 显示所有时钟均为 propagated；
  \item \cmd{check\_timing} 执行默认检查及 \texttt{ideal\_clocks} 检查；
  \item \cmd{report\_clock -skew SD\_DDR\_CLK} 显示 source latency 已计算；
  \item 到 \cmd{sd\_DQ[6]} 的时序报告包含 \texttt{clock network delay (propagated)}，
        slack 由违例变为 \textbf{MET}（1.5143\,ns）。
\end{itemize}
\end{questionBox}

  \item 退出 PrimeTime。
\end{enumerate}

\subsection{任务 5：重新执行运行脚本以减少违例}
\label{lab4:task5}

\noindent\textit{Task 5. Re-Execute the Run Script to Reduce Violation}

\begin{enumerate}
  \item 已提供文件 \cmd{./scripts/orca\_pt\_variables.tcl}，它将完成以下两件事：
\begin{itemize}
  \item 在 \cmd{check\_timing} 默认检查之外，增加会标记 ideal clock 的检查；
  \item 所有创建的时钟均以 propagated clock 方式创建。
\end{itemize}

  \item 从 \cmd{lab4\_clocks} Unix 目录执行运行脚本 \cmd{./RUN.tcl}，
        并将结果记录到日志文件 \cmd{run.log}：
\begin{lstlisting}
unix> pt_shell -f ./RUN.tcl | tee -i run.log
\end{lstlisting}

  \item 启动 PrimeTime 并恢复新保存在 \cmd{./orca\_savesession} 中的会话。

  \item 用适当命令确认下列信息：
\begin{itemize}
  \item \cmd{out\_setup} 违例已减少；
  \item 所有时钟均为 propagated；
  \item 执行 \cmd{check\_timing}，确认除默认检查外还执行 ideal clock 检查；
  \item 现已为 \cmd{SD\_DDR\_CLK} 计算 source latency；
  \item 到 \cmd{sd\_DQ[3]} 的时序报告包含该计算得到的 source latency。
\end{itemize}

\begin{noteBox}
会出现额外违例（更多 setup 违例以及 \cmd{out\_hold} 违例），可忽略。
\end{noteBox}

  \item 退出 PrimeTime。
\end{enumerate}

\noindent 至此 Lab~4 完成。请返回课堂讲义。
